%==============================================================================
\section{Building \& Running the System}
%==============================================================================

This section walks through bringing the whole system up. You can run the dashboard
with real hardware or, for development, with the bundled software simulator.

\subsection{Prerequisites}
\begin{itemize}[noitemsep]
  \item Python 3.10+ on the server machine.
  \item Arduino IDE (or PlatformIO) with the ESP32 board support package, plus the
        \texttt{ArduinoWebsockets} and \texttt{ArduinoJson} libraries -- only needed
        for real hardware.
  \item An ESP32 development board and an AD8232 ECG module with electrodes -- only
        needed for real hardware.
\end{itemize}

\subsection{1. Backend server}
From the \texttt{holter\_monitor/} project directory:
\begin{lstlisting}[style=base,caption={Set up and run the Django/Channels server.}]
# create and activate a virtual environment
python -m venv venv
venv\Scripts\activate            # Windows
# source venv/bin/activate       # macOS / Linux

# install dependencies
pip install django channels daphne numpy scipy PyWavelets websockets

# apply Django's built-in migrations and start the server
python manage.py migrate
python manage.py runserver 0.0.0.0:8000
\end{lstlisting}

Because \texttt{daphne} and \texttt{channels} are installed and
\texttt{ASGI\_APPLICATION} is set, \texttt{runserver} starts the ASGI/WebSocket
server automatically. Binding to \texttt{0.0.0.0} makes the server reachable from
the ESP32 and from other machines on the LAN. Then open
\texttt{http://localhost:8000/} (or \texttt{http://<server-ip>:8000/} from another
device) to load the dashboard.

\subsection{2a. Run with the ESP32 (real ECG)}
\begin{enumerate}[noitemsep]
  \item Open \texttt{arduino/websocket\_ecg\_send\_data/websocket\_ecg\_send\_data.ino}.
  \item Set \texttt{ssid}, \texttt{password} and \texttt{websocket\_server\_host}
        (your server's LAN IP) near the top of the sketch.
  \item Wire the AD8232 output to \texttt{GPIO36} and power the board from 3.3\,V.
  \item Select your ESP32 board and port, then upload.
  \item Open the Serial Monitor at 115200 baud to watch it connect and stream;
        the dashboard trace should begin scrolling.
\end{enumerate}

\subsection{2b. Run without hardware (simulator)}
The repository includes \texttt{holter\_monitor/send\_fake\_data.py}, which connects
to the same WebSocket endpoint and pushes a synthetic ECG-like waveform (a noisy
sine) at the same rate the device would. With the server running:
\begin{lstlisting}[style=base,numbers=none]
python send_fake_data.py
\end{lstlisting}
This is the quickest way to verify the backend, DSP pipeline and dashboard without
any electronics. Note that the simulator feeds the \emph{same} \texttt{ecg\_group},
so the filtered output appears on every open dashboard exactly as real data would.

\subsection{Verifying it works}
\begin{itemize}[noitemsep]
  \item The dashboard page shows a scrolling lime waveform.
  \item The server console logs each received batch and broadcast.
  \item The browser console logs ``WebSocket connected'' and per-batch previews.
  \item With a real AD8232, the QRS spikes of the heartbeat are visible in the
        trace.
\end{itemize}
